\appendix

\chapter{Sensor Datasheet Summary}
\label{app:datasheets}

Table~\ref{tab:app-datasheets} summarizes the key operating specifications of the sensors used in the NirmalNode prototype, for quick reference.

\begin{table}[htbp]
\centering
\caption{Key sensor specifications used in the NirmalNode prototype}
\label{tab:app-datasheets}
\small
\begin{tabular}{@{}p{2.4cm}p{4.8cm}p{2cm}p{2.5cm}p{1.8cm}@{}}
\toprule
\textbf{Sensor} & \textbf{Measured Quantity} & \textbf{Interface} & \textbf{Supply} & \textbf{ESP32 GPIO} \\
\midrule
PMS5003  & PM1.0, PM2.5, PM10 ($\mu$g/m$^3$) + particle counts & UART (3.3~V TTL) & 5~V & GPIO~16/17 \\
MQ135    & NH$_3$, NO$_x$, benzene, CO$_2$ (broad) & Analog (ADC) & 5~V & GPIO~32 \\
MQ136    & H$_2$S                                  & Analog (ADC) & 5~V & GPIO~33 \\
MiCS RED & CO, H$_2$, CH$_4$, NH$_3$              & Analog (ADC) & 3.3--5~V & GPIO~34 \\
MiCS NOX & NO$_2$                                  & Analog (ADC) & 3.3--5~V & GPIO~35 \\
MP135    & General air quality (supplementary)     & Analog (ADC) & 5~V & GPIO~36 \\
Relay    & Fan control (GPIO-driven)               & Digital OUT  & 5~V & GPIO~25 (D25) \\
BME280   & Temperature, humidity, pressure (optional) & I$^2$C    & 3.3~V & SDA/SCL \\
\bottomrule
\end{tabular}
\end{table}

\chapter{ESP32 Firmware Listing}
\label{app:firmware}

The following is an abbreviated listing of the key functions in \texttt{AirGuard\_WiFi.ino}, the WiFi-enabled firmware deployed on the NirmalNode ESP32. The full source code is available in the project repository.

\begin{small}
\begin{verbatim}
// ---- Quick Setup: change only these 3 lines ----
#define WIFI_SSID   "YOUR_WIFI_NAME"
#define WIFI_PASS   "YOUR_WIFI_PASSWORD"
#define SERVER_IP   "192.168.1.100"   // PC IP (ipconfig -> IPv4)
#define RELAY_FAN_PIN  25             // D25 -> Relay IN

// ---- setup() ----
void setup() {
  pinMode(RELAY_FAN_PIN, OUTPUT);
  digitalWrite(RELAY_FAN_PIN, LOW);   // fan OFF at boot
  PMS.begin(9600, SERIAL_8N1, 16, 17);
  analogReadResolution(12);
  connectWiFi();
  calibrateSensors();                 // 50-sample R0 baseline
}

// ---- loop() every 2 seconds ----
void loop() {
  readGasSensors();   // MQ135, MQ136, MiCS(DEMO), MP135
  readPMS5003();      // real UART or DEMO_MODE synthetic PM
  printReport();      // Serial Monitor report
  sendToServer();     // HTTP POST -> ML engine -> fan cmd
  delay(2000);
}

// ---- sendToServer(): POST JSON, receive fan command ----
void sendToServer() {
  StaticJsonDocument<512> doc;
  doc["pm2_5"]  = pmsData.pm2_5;
  doc["mq135_adc"] = mq135RawADC;
  // ... (all 12 sensor fields)
  doc["fan"] = fanOn;
  http.POST(payload);
  // parse response: {"fan":"FAN_ON","aqi":149}
  setFan(strcmp(fanCmd,"FAN_ON") == 0);
}

// ---- setFan(): controls D25 relay ----
void setFan(bool on) {
  fanOn = on;
  digitalWrite(RELAY_FAN_PIN, on ? HIGH : LOW);
}
\end{verbatim}
\end{small}

\chapter{Sample Sensor Record and API Response}
\label{app:sample-data}

The following shows a complete JSON record transmitted from the NirmalNode ESP32 to the server's \texttt{/api/ingest} endpoint during a moderate-pollution event (PM2.5 = 35~$\mu$g/m$^3$, AQI~=~99, purifier active):

\begin{small}
\begin{verbatim}
// --- ESP32 POST payload (every 2 seconds) ---
{
  "timestamp":    "20:16:51",
  "pm1_0":        20,
  "pm2_5":        35,
  "pm10":         51,
  "cnt0_3":       1685,
  "cnt0_5":       490,
  "cnt1_0":       122,
  "cnt2_5":       26,
  "mq135_adc":    1748,
  "mq136_adc":    1170,
  "mics_red_adc": 2219,
  "mics_nox_adc": 1491,
  "mp135_adc":    1624,
  "pms_ok":       true,
  "cal_ok":       true,
  "demo":         true,
  "fan":          false
}

// --- server.py JSON response (same HTTP transaction) ---
{
  "status":  "ok",
  "fan":     "FAN_ON",
  "aqi":     99
}
\end{verbatim}
\end{small}

\noindent The server's \texttt{/api/latest} endpoint simultaneously returns the full ML-enriched telemetry record, including forecast values and online-learning metrics, for consumption by the web dashboard:

\begin{small}
\begin{verbatim}
{
  "timestamp":         "20:16:51",
  "pm2_5":             35,
  "aqi":               99,
  "aqi_category":      "Moderate",
  "aqi_color":         "#f59e0b",
  "forecast_1h":       36.1,
  "forecast_2h":       37.1,
  "forecast_aqi":      102,
  "forecast_aqi_category": "Unhealthy for Sensitive Groups",
  "purifier_trigger":  true,
  "purifier_status":   "ACTIVE (PRE-EMPTIVE)",
  "online_learning": {
    "samples_trained": 144,
    "mae":             11.23,
    "rmse":            18.45,
    "status":          "Model Adapting Live"
  }
}
\end{verbatim}
\end{small}

\chapter{Python ML Engine: AdaptiveAIEngine Class}
\label{app:ml-engine}

The following pseudocode summarizes the core \texttt{process\_stream\_sample()} method of the \texttt{AdaptiveAIEngine} class (\texttt{nirmal\_ml\_engine.py}), which implements the prequential test-then-train protocol:

\begin{small}
\begin{verbatim}
def process_stream_sample(self, sample):
    # 1. Build 14-feature vector (PM + counts + gas ADC + delta + EMA)
    X = self._build_features(sample)

    # 2. Incremental scaler update (online StandardScaler)
    self.scaler.partial_fit(X)
    X_sc = self.scaler.transform(X)

    # 3. PREDICT (if model fitted)
    if self.is_fitted:
        pred = 0.6 * self.sgd.predict(X_sc)
                + 0.4 * self.pa.predict(X_sc)
        blend = min(0.60, (self.n_samples - 10) / 70.0)
        forecast_1h = blend * pred + (1-blend) * self.ema_pm25 * 1.03

    # 4. SCORE (prequential: use previous lag prediction)
    if self.lag_buf:
        past_X_sc, past_pred = self.lag_buf.popleft()
        err = abs(current_pm25 - past_pred)
        self.mae  = cumulative_mae(err)
        self.rmse = cumulative_rmse(err**2)

    # 5. TRAIN (online partial_fit)
    self.sgd.partial_fit(past_X_sc, [current_pm25])
    self.pa.partial_fit(past_X_sc, [current_pm25])

    # 6. AQI + purifier decision
    aqi, category, color = calculate_pm25_aqi(current_pm25)
    purifier = (aqi > 100) or (forecast_1h >= 35.5)

    return { "aqi": aqi, "forecast_1h": forecast_1h,
             "purifier_trigger": purifier, ... }
\end{verbatim}
\end{small}
